\chapter{Fundamentele aplicațiilor web}

\section{Rolurile și Funcțiile Frontend-ului și Backend-ului în Aplicațiile Web}

Frontend-ul și backend-ul reprezintă două componente esențiale ale oricărei aplicații web moderne \cite{frontendBackendWebDev}. Frontend-ul, denumit adesea „partea clientului”, este responsabil pentru interacțiunile utilizatorului și pentru tot ceea ce utilizatorul vede și cu care interacționează direct. Acesta include elementele vizuale, cum ar fi butoane, imagini, formulare și text, precum și funcționalitățile de interacțiune, cum ar fi validarea formularelor și navigarea între pagini. Tehnologiile utilizate frecvent în dezvoltarea frontend-ului includ HTML, CSS, JavaScript și framework-uri moderne precum React, Angular, Vue și Flutter \cite{flutterDev}.

Pe de altă parte, backend-ul, cunoscut și sub numele de „partea serverului”, gestionează logica de afaceri, bazele de date, autentificarea utilizatorilor și comunicarea cu serverele externe. Backend-ul este responsabil pentru procesarea cererilor utilizatorului, accesarea și manipularea datelor stocate, precum și pentru securitatea și gestionarea aplicației. Limbajele comune utilizate pentru dezvoltarea backend-ului includ Node.js, Python, Ruby, Java și Rust \cite{rustProgramming}, fiecare oferind avantaje specifice în funcție de nevoile aplicației.

Această colaborare între frontend și backend este esențială pentru a crea o aplicație web eficientă și scalabilă. Frontend-ul oferă o interfață atractivă și intuitivă, în timp ce backend-ul asigură suportul necesar pentru funcționalitatea corectă și performanța aplicației. Comunicarea dintre cele două componente, de obicei realizată prin API-uri, contribuie la o arhitectură bine definită și la o scalabilitate ușoară a aplicației \cite{frontendBackendWebDev}.



\section{Frontend}
% Aici vei descrie detaliat ce înseamnă frontend-ul unei aplicații. Menționează tehnologii folosite în mod obișnuit, cum ar fi HTML, CSS, JavaScript și framework-uri ca React, Angular, Vue, dar și Flutter, un framework popular pentru dezvoltarea de interfețe mobile și web.

\subsection{Limbaje de programare și tehnologii folosite pentru Frontend}
% Aici vei discuta despre limbajele și tehnologiile comune pentru dezvoltarea frontend-ului, precum HTML, CSS, JavaScript, Flutter și framework-uri moderne ca React și Angular.

\section{Backend}
% Aici vei detalia partea de backend a aplicației, responsabilă de gestionarea logicii afacerii, bazelor de date și middleware. Limbaje precum Python, Node.js, Ruby, Java, și Rust sunt frecvent utilizate.

\subsection{Limbaje de programare și tehnologii folosite pentru Backend}
% Aici vei discuta despre limbajele și framework-urile comune în dezvoltarea backend-ului, precum Node.js, Python (Django/Flask), Ruby on Rails, Java (Spring), și Rust. Rust, cunoscut pentru performanța sa și siguranța memoriei, este din ce în ce mai folosit pentru backend-ul aplicațiilor web.

\subsection{Funcțiile Middleware}
% În această secțiune vei explica rolul middleware-ului în cadrul backend-ului, cum middleware-ul gestionează cererile între client și aplicație și cum ajută la autentificare, logare, validarea datelor și gestionarea erorilor. Exemple de middleware includ autentificarea, gestionarea sesiunilor și optimizarea performanței.

\section{Proiectarea Aplicațiilor}
% În această secțiune vei acoperi principii de proiectare precum SOLID, DRY, KISS și YAGNI, care ajută la menținerea codului curat și ușor de întreținut. De asemenea, poți discuta despre arhitecturi de software cum ar fi MVC, MVP și microservicii. Aceste principii și arhitecturi asigură că aplicațiile web sunt scalabile, ușor de întreținut și că evoluează fără a introduce erori sau probleme de performanță.

\section{Aspecte Cheie ale Aplicațiilor Web}
% În această secțiune finală, vei rezuma principalele aspecte discutate, subliniind importanța unei colaborări eficiente între frontend și backend, precum și folosirea celor mai bune practici de proiectare a aplicațiilor. Discută despre modul în care tehnologiile moderne, cum ar fi Flutter pentru frontend și Rust pentru backend, contribuie la performanța și scalabilitatea aplicațiilor. Încheie prin evidențierea faptului că succesul unei aplicații web depinde de echilibrul între interfața de utilizator (UI/UX) și funcționalitățile de backend bine proiectate.
